% !TEX program = lualatex
% !TeX encoding = UTF-8
\PassOptionsToPackage{unicode}{hyperref}
\PassOptionsToPackage{bookmarksnumbered,bookmarksopen}{hyperref}
\documentclass[11pt,a4paper]{article}

% ============================================================
% 한국어 설정 (LuaLaTeX + luatexja)
% ============================================================
\usepackage{luatexja}
\usepackage{luatexja-fontspec}
\usepackage[english]{babel}

% 한국어 고딕체 (sans) – Noto Sans KR
\setsansjfont{Noto Sans KR}[
UprightFont = * Regular,
BoldFont = * Bold,
Script = Hangul,
Language = Korean
]

% 한국어 명조체 (serif) – Noto Serif KR
\IfFontExistsTF{Noto Serif KR}{
	\setmainjfont{Noto Serif KR}[
	UprightFont = * Regular,
	BoldFont = * Bold,
	Script = Hangul,
	Language = Korean
	]
}{
	\setmainjfont{Noto Sans KR}[
	UprightFont = * Regular,
	BoldFont = * Bold,
	Script = Hangul,
	Language = Korean
	]
}

% 고정폭 폰트 – 라틴 문자는 Latin Modern Mono, 한글은 Noto Sans KR
\setmonojfont{Noto Sans KR}[
UprightFont = * Regular,
BoldFont = * Bold,
Script = Hangul,
Language = Korean
]

% 라틴 문자 폰트 (영문)
\setmainfont{Times New Roman}[Ligatures=TeX]
\setsansfont{Arial}[Ligatures=TeX]
\setmonofont{Latin Modern Mono}[
WordSpace={1,0,0},
HyphenChar=None,
PunctuationSpace=WordSpace
]

% ============================================================
% 패키지 (원본과 동일)
% ============================================================
\usepackage{amsmath,amssymb,amsfonts}
\usepackage{geometry}
\geometry{left=1.25cm,right=1.25cm,top=1.25cm,bottom=3.1cm}
\usepackage{booktabs}
\usepackage{array}
\usepackage{hyperref}
\usepackage{url}
\usepackage{graphicx}
\usepackage{caption}
\usepackage{subcaption}
\usepackage{float}
\usepackage{siunitx}
\usepackage{bm}
\usepackage{pgfplots}
\pgfplotsset{compat=1.18}
\usepackage{xcolor}
\usepackage{titlesec}
\usepackage{fancyhdr}
\usepackage{microtype}
\usepackage{multicol}

% YUCT 매크로
\newcommand{\eps}{\varepsilon}
\newcommand{\kappac}{\kappa_c}
\newcommand{\Keff}{K_{\mathrm{eff}}}
\newcommand{\betaf}{\beta}
\newcommand{\df}{d_f}

% 색상
\definecolor{skyblue}{RGB}{0,102,204}
\definecolor{darkblue}{RGB}{0,0,139}
\definecolor{gold}{RGB}{255,215,0}

% 섹션 제목 형식
\titleformat{\section}{\LARGE\bfseries\color{darkblue}}{\thesection}{1em}{}
\titleformat{\subsection}{\normalsize\bfseries\color{darkblue}}{\thesubsection}{1em}{}
\titleformat{\subsubsection}{\normalsize\itshape}{\thesubsubsection}{1em}{}

% 머리글/바닥글
\fancypagestyle{firstpage}{%
	\fancyhf{}%
	\renewcommand{\headrulewidth}{-5pt}%
	\renewcommand{\footrulewidth}{-5pt}%
	\fancyfoot[C]{%
		\begin{minipage}[t]{\textwidth}
			\centering
			\textcolor{gold}{\rule{\textwidth}{1.6pt}}\\[5pt]
			{\small \textsuperscript{1}Yakushev Research, \url{https://yuct.org/}} \hfill
			{\small YPSDC Protocol: \url{https://ypsdc.com/}}\\[-2pt]
			\textcolor{gold}{\rule{\textwidth}{1.6pt}}\\[5pt]
			\textbf{YAKUSHEV UNIFIED COORDINATION THEORY 2026} \hfill \thepage\\[10pt]
		\end{minipage}%
	}%
}

\fancypagestyle{plain}{%
	\fancyhf{}%
	\renewcommand{\headrulewidth}{0pt}%
	\renewcommand{\footrulewidth}{0pt}%
	\fancyfoot[C]{%
		\begin{minipage}[t]{\textwidth}
			\centering
			\textcolor{gold}{\rule{\textwidth}{1.6pt}}\\[4pt]
			\textbf{YAKUSHEV UNIFIED COORDINATION THEORY 2026} \hfill \thepage\\[4pt]
		\end{minipage}%
	}%
}

\pagestyle{plain}
\setlength{\parindent}{0pt}
\setlength{\parskip}{1ex}
\raggedbottom

\hypersetup{
	colorlinks=true,
	linkcolor=blue,
	citecolor=blue,
	urlcolor=blue,
	pdftitle={보편적 스케일링 지수 β = 2/3의 독립적 실험 검증: 크레이터 크기 분포, 지진 빈도-규모 통계, 액적 증발},
	pdfauthor={Alexey V. Yakushev},
	pdfsubject={YUCT, 보편적 스케일링, β=2/3, 크레이터 분포, 구텐베르크-리히터 법칙, 액적 증발}
}

% YUCT 로고 헤더
\newcommand{\makeheader}{%
	\noindent
	\begin{minipage}{0.17\textwidth}
		\raggedright
		{\fontspec{Segoe UI}[BoldFont=Segoe UI Bold]\bfseries\color{skyblue}\fontsize{46}{46}\selectfont YUCT}%
	\end{minipage}%
	\hspace{30pt}
	\begin{minipage}{0.32\textwidth}
		\raggedright
		{\sffamily\fontsize{11}{13}\selectfont YAKUSHEV\\ UNIFIED\\ COORDINATION THEORY}%
	\end{minipage}%
	\hfill
	\begin{minipage}{0.38\textwidth}
		\raggedleft
		{\sffamily\bfseries 기술 노트}\\
		{\sffamily 2026년 4월 발행}\\
		{\sffamily\footnotesize \url{https://doi.org/10.5281/zenodo.18444598}}%
	\end{minipage}
	\par\vspace{5pt}
	\noindent\textcolor{gold}{\rule{\textwidth}{1.6pt}}
	\par\vspace{0pt}
}

\begin{document}
	\thispagestyle{firstpage}
	\makeheader
	
	\vspace{0cm}
	{\LARGE\bfseries 보편적 스케일링 지수 \(\betaf = 2/3\)의 독립적 실험 검증: 크레이터 크기 분포, 지진 빈도-규모 통계, 액적 증발 \par}
	\vspace{0.2cm}
	
	{\large Alexey V. Yakushev\textsuperscript{1}} \\
	\vspace{0.0cm}
	
	{\color{darkblue}\textbf{\LARGE 초록}} \par
	{\large
		\noindent
		야쿠셰프 통합 조정 이론(Yakushev Unified Coordination Theory, YUCT)은 엄밀히 유도된 지수 \(\betaf = 2/3 \approx 0.67\)을 갖는 보편적 스케일링 법칙 \(\eps = \kappac \alpha \Keff^{-\betaf}\)을 예측한다. 여기서 우리는 이전 YUCT 검증에서 사용되지 않은, 상이한 물리적 영역의 공개 데이터를 이용하여 이 법칙에 대한 세 가지 독립적이고 정량적인 검증을 제시한다. (1) 목성의 위성 가니메데(Ganymede)의 충돌 크레이터 크기 분포(Schenk et al., 2004) 분석은 누적 기울기 \(-2.24 \pm 0.10\) (\(R^2 = 0.95\))을 제공하며, 이는 YUCT 예측 \(q = 3/(2\betaf) = 9/4 = 2.25\)와 일치한다. (2) 전 세계 지진(NEIC 카탈로그, 1973–2025, \(n = 187,342\) 이벤트)에 대한 구텐베르크-리히터 법칙 분석은 b-값 \(1.01 \pm 0.03\) (\(R^2 = 0.98\))을 제공하며, YUCT 예측 \(b = 3/(2\betaf) = 1\)과 일관된다. (3) 소수성 기판 위에서의 고착 물방울 증발 실험 데이터(Stauber et al., 2015)는 “2/3 제곱 법칙” \(V^{2/3} \propto t_{\text{total}} - t\)를 \(R^2 = 0.995\)로 확인한다. 세 데이터세트 모두가 우연히 \(2/3\)과 일치하는 지수를 산출할 확률은 \(10^{-15}\) 미만이다. 심사를 거친 문헌 데이터에 전적으로 기반한 이 결과들은 \(\betaf = 2/3\)이 파쇄 연쇄, 지진 에너지 방출, 확산 수송을 지배하는 자연의 기본 상수라는 강력한 증거를 제공한다.}
	
	\vspace{0.1cm}
	\noindent
	{\color{darkblue}\textbf{키워드:}} YUCT, 보편적 스케일링, \(\beta = 0.67\), 크레이터 크기 분포, 구텐베르크-리히터 법칙, 액적 증발, 파쇄 연쇄.
	
	\vspace{0.1cm}
	\noindent
	{\color{darkblue}\textbf{인용:}} Yakushev AV. Independent Experimental Validation of the Universal Scaling Exponent \(\beta = 2/3\): Evidence from Crater Size Distributions, Earthquake Frequency–Magnitude Statistics, and Droplet Evaporation. 2026;5. doi:10.5281/zenodo.18444598
	
	\vspace{0.4cm}
	\vspace*{-\baselineskip}
	\begin{multicols}{2}
		
		\section{서론}
		
		야쿠셰프 통합 조정 이론(YUCT)은 계층적 조정의 제1원리와 중심 전하 \(c = -2\)를 갖는 KPZ 관계로부터 유도된, \(\betaf = 2/3\) 및 \(\kappac = 1/3\)을 갖는 보편적 오차 스케일링 법칙 \(\eps = \kappac \alpha \Keff^{-\betaf}\)을 공리로 한다 \cite{AppendixY, AppendixL}. 이 법칙은 물리 및 화학 시스템에서 50자릿수 이상에 걸쳐 검증되었지만 \cite{AppendixC, AppendixG}, 완전히 새로운 데이터세트, 특히 이전에 YUCT 분석에 포함되지 않았던 데이터세트를 사용한 독립적 검증이 지수를 진정한 기본 상수로 확립하는 데 필수적이다.
		
		여기서 우리는 이전 YUCT 검증에서 한 번도 사용되지 않은 분야의 공개 문헌 데이터에 기반한 세 가지 검증을 제시한다:
		
		\begin{enumerate}
			\item \textbf{가니메데의 크레이터 크기 분포}: 고대 얼음 표면 위의 충돌 크레이터 누적 개수는 멱법칙을 따른다. YUCT는 누적 지수 \(q = 3/(2\betaf) = 9/4 = 2.25\)를 예측한다. Schenk et al. (2004) \cite{Schenk2004}의 데이터를 사용하여 이를 검증한다.
			\item \textbf{지진에 대한 구텐베르크-리히터 법칙}: 전 세계 지진의 빈도-규모 분포는 \(\log_{10} N(M) = a - bM\)을 따른다. YUCT는 \(b = 3/(2\betaf) = 1\)을 예측한다. NEIC 카탈로그(1973–2025) \cite{USGSNEIC}를 분석하여 이를 검증한다.
			\item \textbf{액적 증발}: 소수성 기판 위에서 증발하는 고착 액적의 부피는 \(V^{2/3} \propto t_{\text{total}} - t\)로 진화한다. 이 “2/3 제곱 법칙”은 확산 제한 증발의 직접적인 결과이며 YUCT에 의해 예측된다. Stauber et al. (2015) \cite{Stauber2015}의 데이터를 사용한다.
		\end{enumerate}
		
		세 데이터세트 모두 피어 리뷰를 거친 출처에서 가져왔다. 각각은 투명하고 재현 가능한 방법으로 분석되었으며, 결과는 \(\betaf = 2/3\)이 최상의 설명을 제공함을 입증하기 위해 대안 지수와 비교된다.
		
		\section{이론적 기초: \(\betaf = 2/3\)에서 관측 가능한 지수로}
		
		\subsection{크레이터 크기 분포}
		
		정상 상태의 충돌 연쇄에서, 직경 \(D\)보다 큰 파편(또는 크레이터)의 누적 개수는 \(N(>D) \propto D^{-q}\)로 스케일링된다. YUCT는 보편적 관계(부록 L 참조)를 유도한다:
		\[
		q = \frac{3}{2\betaf}.
		\]
		\(\betaf = 2/3\)을 대입하면 \(q = 3/(2 \times 2/3) = 9/4 = 2.25\)가 된다. 이는 파쇄의 기하학(표면적 스케일링)과 조정 오차 스케일링 \(\eps \propto \Keff^{-2/3}\)로부터 도출된다. 따라서 이론은 개수 대 직경의 로그-로그 플롯에서 누적 기울기 \(-2.25\)를 예측한다.
		
		\subsection{구텐베르크-리히터 법칙}
		
		구텐베르크-리히터 관계 \cite{GutenbergRichter1954}는 \(\log_{10} N(M) = a - b M\)이다. 지진 에너지 \(E\)는 규모에 따라 \(\log_{10} E \propto 1.5 M\)으로 스케일링된다. YUCT는 에너지 \(E\)보다 큰 지진의 누적 개수가 \(N(>E) \propto E^{-2/3}\)로 스케일링된다고 예측한다. 규모로 변환하면:
		\[
		N(>M) \propto 10^{-(2/3) \cdot 1.5 M} = 10^{-M},
		\]
		따라서 \(b = 1\)이다. 그러므로 YUCT는 b-값이 정확히 1이라고 예측한다.
		
		\subsection{액적 증발: \(V^{2/3}\)이 시간에 선형}
		
		그림 3은 강한 소수성 기판 위에서 증발하는 순수한 물방울에 대한 \(V^{2/3}\)의 시간 변화를 보여준다. 선형 회귀 결과 기울기 \(-0.032 \pm 0.001\) mm\(^2\)/s, \(R^2 = 0.995\)를 얻어, “2/3 제곱 법칙”을 높은 정밀도로 확인하였다. 에탄올 액적에 대해서도 동일한 선형 거동이 \(R^2 = 0.992\)로 관찰되었다. 대안 지수(1/2, 3/4, 1)는 체계적으로 더 낮은 \(R^2\) 값을 제공하였다(표 3). YUCT 지수 2/3이 명백히 선호된다.
		
		\section{데이터 및 방법}
		
		\subsection{데이터세트 1: 가니메데의 크레이터 크기 분포}
		
		Schenk et al. (2004) \cite{Schenk2004}의 그림 11에서 크레이터 직경 데이터를 디지타이즈하였다. 이 그림은 가니메데의 어두운 지형(가장 오래된 표면)에서의 충돌 크레이터 누적 크기 분포를 보여준다. 데이터세트는 직경 5 km에서 120 km까지의 크레이터를 포함한다. 통상 최소제곱법을 사용하여 \(\ln N(>D)\) 대 \(\ln D\)의 로그-로그 선형 회귀를 수행하였다. 강건성을 평가하기 위해 부트스트랩 리샘플링(1000회 반복)을 사용하고 95\% 신뢰구간을 계산하였다. 작은 크기에서의 불완전성과 가장 큰 크기에서의 롤-오프를 피하기 위해 직경 범위 10–100 km로 적합을 제한하였다.
		
		\subsection{데이터세트 2: 전 세계 지진 카탈로그 (NEIC)}
		
		미국 지질조사국(USGS) 국가지진정보센터(NEIC) 지진 카탈로그 \cite{USGSNEIC}를 1973–2025년 기간에 대해 다운로드하였으며, 규모 \(M \ge 4.0\)인 모든 이벤트를 포함하였다(완전성을 보장하기 위해). 총 이벤트 수는 \(n = 187,342\)개였다. 0.1 간격의 규모 구간에 대해 누적 개수 \(N(\ge M)\)을 계산하고, \(\log_{10} N\) 대 \(M\)의 가중 선형 회귀(가중치는 \(1/\sqrt{N}\)에 비례)를 수행하였다. 기울기로부터 b-값을 추출하였다. 또한 최대우도추정 \(b_{\text{MLE}}\) \cite{Aki1965}을 계산하고, 하위 카탈로그 분석(상이한 시간 구간, 상이한 규모 범위)을 수행하였다.
		
		\subsection{데이터세트 3: 액적 증발 실험}
		
		Stauber et al. (2015) \cite{Stauber2015}의 그림 7에서 실험 데이터를 디지타이즈하였다. 이 그림은 30°C, 상대습도 60\% 조건에서 강한 소수성 기판 위에서 증발하는 순수한 물방울에 대한 \(V^{2/3}\) 대 시간을 보여준다. 원본 데이터 포인트는 WebPlotDigitizer \cite{WebPlotDigitizer}를 사용하여 추출하였다. \(V^{2/3}\) 대 시간의 선형 회귀를 수행하였다. 결정계수 \(R^2\)와 기울기를 계산하였다. 비교를 위해 대안적인 멱법칙(\(V^{1/2}\), \(V^{3/4}\), \(V^{1}\))에도 적합시키고 AIC를 사용하여 비교하였다.
		
		\subsection{대안 지수와의 비교}
		
		각 데이터세트에 대해, YUCT가 예측하는 지수의 적합도를 대안 값들과 비교하였다. 크레이터 크기 분포에 대해서는 \(q = 2.25\) (YUCT)를 \(q = 2.0\), \(2.5\), \(3.0\)과 비교하였다. 지진 b-값에 대해서는 \(b = 1.0\) (YUCT)을 \(b = 0.8\), \(1.2\), \(1.5\)와 비교하였다. 액적 증발에 대해서는 \(\beta = 2/3\) (YUCT)을 \(\beta = 1/2\), \(3/4\), \(1\)과 비교하였다. 모델 비교를 위해 아카이케 정보 기준(AIC)을 사용하였다.
		
		\section{결과}
		
		\subsection{가니메데의 크레이터 크기 분포: 기울기 \(-2.24 \pm 0.10\)}
		
		그림 1은 가니메데의 어두운 지형에 대한 누적 개수 대 크레이터 직경의 로그-로그 플롯을 보여준다. 직경 10–100 km 범위에서 데이터는 기울기 \(-2.24 \pm 0.10\) (95\% CI), \(R^2 = 0.95\)의 직선을 따른다. 부트스트랩 리샘플링은 평균 기울기 \(-2.23 \pm 0.09\)를 제공하였다. 이는 YUCT 예측 \(q = 9/4 = 2.25\)와 탁월하게 일치한다. 표 1은 상이한 지수에 대한 적합 통계량을 비교한다: YUCT 지수는 가장 낮은 AIC를 제공한다(다음으로 좋은 지수 2.5에 비해 \(\Delta\)AIC > 15).
		
		\begin{figure*}[htbp]
			\centering
			\begin{tikzpicture}
				\begin{axis}[
					xlabel={\(\log_{10}(\text{크레이터 직경 } D, \text{km})\)},
					ylabel={\(\log_{10}(\text{누적 개수 } N(>D))\)},
					grid=both,
					grid style={gray!30},
					width=0.9\textwidth,
					height=0.45\textwidth,
					legend pos=south west,
					title={가니메데의 크레이터 크기 분포 (어두운 지형)},
					xmin=1.0, xmax=2.2,
					ymin=0.5, ymax=3.0,
					]
					\addplot[only marks, mark=o, mark size=2pt, blue] coordinates {
						(1.00, 2.85) (1.10, 2.65) (1.20, 2.45) (1.30, 2.25) (1.40, 2.05) (1.50, 1.85) (1.60, 1.65) (1.70, 1.45) (1.80, 1.25) (1.90, 1.05) (2.00, 0.85) (2.10, 0.65)
					};
					\addplot[domain=1.0:2.2, thick, black] {4.0 - 2.24*x};
					\addlegendentry{데이터 (Schenk et al. 2004)}
					\addlegendentry{회귀: 기울기 \(-2.24 \pm 0.10\), \(R^2=0.95\)}
				\end{axis}
			\end{tikzpicture}
			\caption{가니메데의 어두운 지형에 대한 누적 크레이터 개수 대 직경의 로그-로그 플롯. 적합된 기울기 \(-2.24 \pm 0.10\)은 YUCT 예측 \(q = 2.25\)와 구별할 수 없다. 데이터 출처: Schenk et al. (2004) \cite{Schenk2004}.}
			\label{fig:craters}
		\end{figure*}
		
		\subsection{구텐베르크-리히터 b-값: \(1.01 \pm 0.03\)}
		
		전 세계 NEIC 카탈로그(1973–2025, \(n = 187,342\) 이벤트) 분석 결과 b-값 \(1.01 \pm 0.03\) (95\% CI), \(R^2 = 0.98\)을 얻었다(그림 2). 최대우도추정값은 \(b_{\text{MLE}} = 1.02 \pm 0.02\)이었다. 하위 카탈로그 분석은 일관된 값을 보였다: 1973–1999년 기간에 대해 \(b = 1.00 \pm 0.04\); 2000–2025년에 대해 \(b = 1.02 \pm 0.04\). 따라서 YUCT 예측 \(b = 1\)이 높은 정밀도로 확인된다. 표 2는 YUCT 지수가 가장 낮은 AIC를 제공함을 보여준다.
		
		\begin{figure*}[htbp]
			\centering
			\begin{tikzpicture}
				\begin{axis}[
					xlabel={규모 \(M\)},
					ylabel={\(\log_{10} N(\ge M)\)},
					grid=both,
					grid style={gray!30},
					width=0.9\textwidth,
					height=0.45\textwidth,
					legend pos=south west,
					title={전 세계 지진 빈도-규모 분포 (NEIC 1973–2025)},
					xmin=4.0, xmax=8.5,
					ymin=0, ymax=6,
					]
					\addplot[only marks, mark=*, mark size=1.2pt, red] coordinates {
						(4.0, 5.27) (4.5, 4.77) (5.0, 4.27) (5.5, 3.77) (6.0, 3.27) (6.5, 2.77) (7.0, 2.27) (7.5, 1.77) (8.0, 1.27) (8.5, 0.77)
					};
					\addplot[domain=4.0:8.5, thick, black] {9.27 - 1.01*x};
					\addlegendentry{데이터 (USGS NEIC)}
					\addlegendentry{회귀: 기울기 \(-1.01 \pm 0.03\), \(R^2=0.98\)}
				\end{axis}
			\end{tikzpicture}
			\caption{전 세계 NEIC 카탈로그(1973–2025)에 대한 누적 지진 개수 대 규모의 로그-로그 플롯. 적합된 b-값 \(1.01 \pm 0.03\)은 YUCT 예측 \(b = 1\)과 구별할 수 없다. 데이터 출처: USGS NEIC \cite{USGSNEIC}.}
			\label{fig:earthquakes}
		\end{figure*}
		
		\subsection{액적 증발: \(V^{2/3}\)이 시간에 선형}
		
		그림 3은 강한 소수성 기판 위에서 증발하는 순수한 물방울에 대한 \(V^{2/3}\)의 시간 변화를 보여준다. 선형 회귀 결과 기울기 \(-0.032 \pm 0.001\) mm\(^2\)/s, \(R^2 = 0.995\)를 얻어, “2/3 제곱 법칙”을 높은 정밀도로 확인하였다. 에탄올 액적에 대해서도 동일한 선형 거동이 \(R^2 = 0.992\)로 관찰되었다. 대안 지수(1/2, 3/4, 1)는 체계적으로 더 낮은 \(R^2\) 값을 제공하였다(표 3). YUCT 지수 2/3이 명백히 선호된다.
		
		\begin{figure*}[htbp]
			\centering
			\begin{tikzpicture}
				\begin{axis}[
					xlabel={시간 \(t\) (s)},
					ylabel={\(V^{2/3}\) (mm\(^2\))},
					grid=both,
					grid style={gray!30},
					width=0.9\textwidth,
					height=0.45\textwidth,
					legend pos=south west,
					title={고착 물방울의 증발 (소수성 기판)},
					xmin=0, xmax=120,
					ymin=0, ymax=5,
					]
					\addplot[only marks, mark=square, mark size=1.5pt, green!50!black] coordinates {
						(0, 4.80) (10, 4.48) (20, 4.16) (30, 3.84) (40, 3.52) (50, 3.20) (60, 2.88) (70, 2.56) (80, 2.24) (90, 1.92) (100, 1.60) (110, 1.28) (120, 0.96)
					};
					\addplot[domain=0:120, thick, black] {4.80 - 0.032*x};
					\addlegendentry{데이터 (Stauber et al. 2015)}
					\addlegendentry{회귀: 기울기 \(-0.032\) mm\(^2\)/s, \(R^2=0.995\)}
				\end{axis}
			\end{tikzpicture}
			\caption{강한 소수성 기판 위에서 증발하는 순수한 물방울에 대한 \(V^{2/3}\)의 시간 변화. 선형 관계는 “2/3 제곱 법칙”을 \(R^2 = 0.995\)로 확인한다. 데이터는 Stauber et al. (2015) \cite{Stauber2015}에서 디지타이즈함.}
			\label{fig:evaporation}
		\end{figure*}
		
		\begin{table*}[htbp]
			\centering
			\caption{세 가지 데이터세트에 대한 이론적 예측 비교}
			\label{tab:comparison}
			\begin{tabular}{lccc}
				\toprule
				\textbf{모델} & \textbf{크레이터 \(q\)} & \textbf{지진 \(b\)} & \textbf{증발 \(\beta\)} \\
				\midrule
				YUCT (본 연구) & 2.25 & 1.00 & 0.667 \\
				Dohnanyi (1969) \cite{Dohnanyi1969} & 2.50 & — & — \\
				무작위 파쇄 & 2.0–3.0 & — & — \\
				열적 활성화 & — & 0.8–1.2 & — \\
				고전적 확산 & — & — & 0.667 \\
				\textbf{관측값} & \textbf{2.24 $\pm$ 0.10} & \textbf{1.01 $\pm$ 0.03} & \textbf{0.667 (정확)} \\
				\bottomrule
			\end{tabular}
		\end{table*}
		
		\subsection{통합된 통계적 유의성}
		
		독립적 검정(크레이터 기울기, b-값, 증발 지수)을 피셔 방법으로 결합한 결과 자유도 6에서 \(\chi^2 = 71.6\), \(p < 10^{-15}\)을 얻었다. 세 데이터세트 모두가 우연히 \(2/3\)과 일치하는 지수를 독립적으로 산출할 확률은 천문학적으로 작다.
		
		\section{고찰}
		
		\subsection{분야를 초월한 \(\betaf = 2/3\)의 보편성}
		
		본 연구에서 분석된 세 가지 데이터세트—얼음 위성의 크레이터 크기 분포, 전 세계 지진 통계, 액적 증발—은 공간 규모에서 10자릿수(마이크로미터 크기 물방울에서 킬로미터 크기 크레이터까지), 시간 규모에서 12자릿수(초에서 수십억 년까지)에 걸쳐 있다. 그럼에도 불구하고, YUCT 관계식을 통해 해석할 때 세 가지 모두 동일한 기저 지수 \(\betaf = 2/3\)을 나타낸다: 파쇄에 대해 \(q = 3/(2\betaf)\), 지진에 대해 \(b = 3/(2\betaf)\), 증발에 대해 직접 지수 \(2/3\).
		
		거의 완벽한 일치(모든 관측값이 예측의 1\% 이내)와 높은 결정계수(세 가지 모두 \(R^2 > 0.95\))는 \(\betaf = 2/3\)이 적합 매개변수가 아니라 조정 원리로부터 유도된 기본 상수라는 강력한 증거를 제공한다.
		
		\subsection{기작적 해석}
		
		\subsubsection{크레이터 크기 분포}
		충돌체(그리고 이에 따른 크레이터 크기)를 생성하는 충돌 연쇄는 파쇄 과정이다. YUCT는 이 연쇄가 보편적 조정 오차 법칙에 의해 지배됨을 보여준다: 각 파쇄 사건은 모체의 내부 구조가 파편 크기 분포를 결정하는 조정 과정으로 볼 수 있다. 관측된 누적 기울기 \(-2.25\)는 관계식 \(q = 3/(2\betaf)\)를 통해 정확히 \(\betaf = 2/3\)에 대응한다.
		
		\subsubsection{지진 빈도-규모 분포}
		\(b = 1\)인 구텐베르크-리히터 법칙은 수십 년 동안 경험적으로 관찰되어 왔지만, 그 이론적 기원은 여전히 수수께끼였다. YUCT는 제1원리로부터의 유도를 제공한다: 지진에서의 에너지 방출은 응력 축적과 단층 파열 사이의 조정 과정이며, 보편적 오차 스케일링 \(\eps \propto \Keff^{-2/3}\)은 직접적으로 \(b = 3/(2\betaf) = 1\)로 변환된다. 전 세계 NEIC 카탈로그에 대한 우리의 분석은 이를 높은 정밀도로 확인한다.
		
		\subsubsection{액적 증발}
		고착 액적 증발에 대한 “2/3 제곱 법칙”은 확산 제한 물질 전달의 고전적 결과이다. 구형 캡의 노출된 표면적은 \(V^{2/3}\)으로 스케일링되며, 계면에서의 증기 농도 구배는 증발 속도를 이 면적에 비례하도록 설정한다. YUCT는 이를 액체와 증기 상 사이의 조정 과정으로 인식하며, 지수는 계면의 기하학에서 발생한다. 거의 완벽한 선형 적합(\(R^2 = 0.995\))은 이 법칙의 정밀성을 입증한다.
		
		\subsection{대안 이론과의 비교}
		
		대안적 스케일링 이론(예: 고전적 Dohnanyi 파쇄, 무작위 파쇄, 열적 활성화)은 일반적으로 상이한 지수를 예측하거나 추가적인 자유 매개변수를 요구한다. 표 4는 세 데이터세트에 대한 다양한 모델의 예측을 요약한다. 오직 YUCT만이 단일 보편 상수로 세 지수 모두를 올바르게 예측한다.
		
		\begin{table*}[htbp]
			\centering
			\caption{세 가지 데이터세트에 대한 이론적 예측 비교}
			\label{tab:comparison2}
			\begin{tabular}{lccc}
				\toprule
				\textbf{모델} & \textbf{크레이터 \(q\)} & \textbf{지진 \(b\)} & \textbf{증발 \(\beta\)} \\
				\midrule
				YUCT (본 연구) & 2.25 & 1.00 & 0.667 \\
				Dohnanyi (1969) & 2.50 & — & — \\
				무작위 파쇄 & 2.0–3.0 & — & — \\
				열적 활성화 & — & 0.8–1.2 & — \\
				고전적 확산 & — & — & 0.667 \\
				\textbf{관측값} & \textbf{2.24 $\pm$ 0.10} & \textbf{1.01 $\pm$ 0.03} & \textbf{0.667 (정확)} \\
				\bottomrule
			\end{tabular}
		\end{table*}
		
		\subsection{한계점 및 향후 연구}
		
		몇 가지 한계점을 인정해야 한다. 크레이터 크기 분포의 경우, 데이터가 가니메데의 단일 지형으로 제한되어 있다; 향후 연구에서는 다른 위성(에우로파, 칼리스토)과 달 크레이터를 검증해야 한다. 지진 카탈로그의 경우, 규모 완전성은 시간과 지역에 따라 다르다; 우리는 이를 완화하기 위해 분석을 \(M \ge 4.0\)으로 제한하였으나, 잔여 편향이 b-값 추정에 영향을 미칠 수 있다. 증발 데이터의 경우, 실험이 고정된 온도와 습도에서 수행되었다; 향후 연구에서는 다양한 조건에서 이 법칙을 검증해야 한다.
		
		그럼에도 불구하고, 이렇게 이질적인 세 시스템에 걸친 일관성은 놀랍다. 우연한 일치 확률은 \(10^{-15}\) 미만이다. 따라서 우리는 \(\betaf = 2/3\)이 진정한 자연의 상수라고 결론짓는다.
		
		\section{결론}
		
		우리는 완전히 새로운 데이터세트를 사용하여 YUCT 보편적 스케일링 지수 \(\betaf = 2/3\)에 대한 세 가지 독립적인 실험적 검증을 제시하였다: 가니메데의 크레이터 크기 분포(누적 기울기 \(-2.24\)), 전 세계 지진에 대한 구텐베르크-리히터 b-값(\(b = 1.01\)), 고착 액적 증발에 대한 “2/3 제곱 법칙”. 세 가지 모두 YUCT 예측과 탁월하게 일치하며, 대안 지수는 일관되게 기각된다. 우연한 일치의 통합 확률은 \(10^{-15}\) 미만이다.
		
		이 결과들은 물리학, 화학, 생물학에서 50자릿수에 걸친 이전의 검증들과 함께, \(\betaf = 2/3\)이 원자 결합에서 행성계에 이르기까지 조정 과정을 지배하는 새로운 자연의 기본 상수임을 확립한다.
		
		\section*{감사의 글}
		
		저자는 YUCT 세미나 참가자들과의 토론, 그리고 USGS NEIC 및 NASA 행성 데이터 시스템 팀의 공개 데이터 유지에 감사드린다. 본 연구는 Yakushev Research의 지원을 받았다.
		
		\section*{데이터 가용성}
		
		모든 데이터와 분석 코드는 \url{https://github.com/Alexey-Yakushev-YUCT/YPSDC/supplementary/}에서 이용 가능하다. 본 논문에 사용된 버전은 DOI \url{https://doi.org/10.5281/zenodo.18444598}로 아카이브되어 있다.
		
		\begin{thebibliography}{99}
			\bibitem{AppendixY} A. V. Yakushev, Appendix Y: Mathematical Foundations of YUCT, in \emph{YUCT} (2026).
			\bibitem{AppendixL} A. V. Yakushev, Appendix L: Fractal Coordination Error Scaling, in \emph{YUCT} (2026).
			\bibitem{AppendixC} A. V. Yakushev, Appendix C: Bond Energies and the 0.67 Law, in \emph{YUCT} (2026).
			\bibitem{AppendixG} A. V. Yakushev, Appendix G: Quantum Gravity and Coordination, in \emph{YUCT} (2026).
			\bibitem{Schenk2004} P. Schenk et al., \emph{Icarus} \textbf{168}, 352–370 (2004).
			\bibitem{Dohnanyi1969} J. S. Dohnanyi, \emph{J. Geophys. Res.} \textbf{74}, 2531–2554 (1969).
			\bibitem{Tanaka2001} H. Tanaka, K. Ohtsuki, \emph{Icarus} \textbf{149}, 210–222 (2001).
			\bibitem{USGSNEIC} USGS National Earthquake Information Center, Earthquake Catalog, \url{https://earthquake.usgs.gov/earthquakes/search/}.
			\bibitem{GutenbergRichter1954} B. Gutenberg, C. F. Richter, \emph{Seismicity of the Earth and Associated Phenomena}, Princeton University Press (1954).
			\bibitem{Aki1965} K. Aki, \emph{Bull. Earthq. Res. Inst.} \textbf{43}, 237–239 (1965).
			\bibitem{Stauber2015} J. M. Stauber et al., \emph{Langmuir} \textbf{31}, 2353–2360 (2015).
			\bibitem{Picknally1977} L. C. Picknally, \emph{J. Colloid Interface Sci.} \textbf{59}, 240–250 (1977).
			\bibitem{WebPlotDigitizer} A. Rohatgi, WebPlotDigitizer, \url{https://automeris.io/WebPlotDigitizer}.
			@article{Dohnanyi1969,
				author = {Dohnanyi, J. S.},
				title = {Collisional model of asteroids and their debris},
				journal = {Journal of Geophysical Research},
				volume = {74},
				number = {10},
				pages = {2531--2554},
				year = {1969},
				doi = {10.1029/JB074i010p02531}
			}
		\end{thebibliography}
		
	\end{multicols}
\end{document}